\section{Discussion}\label{sec:discussion}

The memory accounting of Sec.~\ref{sec:algorithm} holds in practice. With an implicit model and no distance map, the working set is the path stack, the in-edge buffers, the divided-difference workspace and the output accumulators, a few kilobytes on top of a 4~MB process, and it does not change between $N=10^4$ and $N=2^{64}$. The state vector is never formed, and neither is a Krylov basis. Three things are still stored and should not be forgotten: an explicit model, when the model is explicit; the distance map, when one is used, which holds the states of a ball of radius $r_d$ and cost up to about 1.5~MB at $r_d=3$ and 3~MB at $r_d=4$ here; and the outputs, which scale with the number of sensors, inputs, patterns and time samples. The memory figures quoted in this paper include all three.

Whether the method is useful is decided by the number of walk prefixes that have to be visited to reach a certificate, and that number depends on four things. The first is the comparison ratio $\mu$, which sets how fast the contribution of long walks decays. The cost grows steeply as $\mu$ approaches one, and for $\mu\ge1$ only the finite-horizon bound is available, with a depth that grows with $\beta t$. On the directed chains of Sec.~\ref{sec:hard}, where back-coupling creates many excursions, the walk sum was slower than \texttt{expm\_multiply} on the whole chain by two to three orders of magnitude or did not certify at all. The second is the in-degree, which sets the branching. The third is the distance information available to the bounds: on the networks a radius-3 map reduced the time to a certificate by factors of about 10 to several hundred, for at most about 1.5~MB. The fourth, which decides the comparison with local methods, is how often distinct walks end at the same state. A local state vector merges all walks that reach a state; walk enumeration cannot, because the temporal factor of a walk depends on its whole multiset of diagonal entries. On a lattice this costs the walk sum orders of magnitude. On a locally tree-like graph few walks merge, the set of visited prefixes and the adaptive set are of similar size, and the walk sum then gains from doing temporal work only where a prefix reaches a source. A local method, in contrast, integrates every retained state over the whole time grid. The uniform ball is the weakest of the three on such graphs because its size is set by its radius, $(\Delta-1)^R$, rather than by the weights. The runtime advantage over the adaptive set also held in the tested degree-16 transfer responses between separated states; the degree-4 cases were too fast to rank reliably (Sec.~\ref{sec:noncol}).

The memory comparison is less one-sided than the extreme cases suggest. The walk sum's memory stays nearly fixed, between 4.5 and 6.1~MB here including the distance map, while the local methods' memory grows with their sets; tiny local sets are in fact slightly leaner than the walker, and the large savings (a factor of 9 against the adaptive set and of about 500 against the ball at degree 32, $\mu=0.57$, target $10^{-4}$, and 3 to 29 against the adaptive set over the realizations of Table~\ref{tab:ensemble}) appear only once the retained sets reach $10^5$ to $10^7$ states.

Several limitations are worth stating plainly. The certificates are worst-case over the endpoints of the discarded walks. Against \texttt{expm\_multiply} references (Fig.~\ref{fig:certificate}; no distance map on the networks) they were loose by one to two orders of magnitude on the lattice and by up to four on the tree-like networks, where a discarded walk rarely returns to the sensor but is charged as if it did. The comparisons with tighter walk sums in Table~\ref{tab:ensemble} provide upper bounds on certificate looseness, up to $2\times10^6$ at degree 32, but do not determine its magnitude. Because the cost grows faster than the inverse certificate, tighter bounds would translate directly into lower cost, and the largest gain is to be had on exactly the graphs where the walk sum is competitive. Any improved completion bound applies equally to the adaptive set, whose certificate (Proposition~\ref{prop:local}) uses the same completion estimate, so whether it would favor walk enumeration is an open question. The distance oracle already enters the certificate, through the factor $\mu^{\delta_k}$ for the shortest completion, but beyond that the bound lets every step of a completion head back toward the sources. Bounds that use how few completions actually return, for instance through return probabilities on locally tree-like graphs, are the natural next step. The uniform-in-time bound needs $\mu<1$; diagonal scaling can extend it, but an optimized scaling vector is itself a state-sized object. Oscillatory systems produce cancellation between walks, which we monitor but do not remove. The scalar diagonal splitting integrates only scalar local dynamics exactly; second-order units were handled here by a local modal transformation, and a genuine block-diagonal splitting would require divided differences of matrix arguments. Finally, the method computes responses at prescribed times; it does not by itself provide controllability measures, reduced models or closed-loop behavior.

Two extensions are immediate from the formulas but were not developed here. Parametric sensitivities come almost for free: a walk visiting state $i$ $m_i$ times satisfies $\partial_{d_i}f_t[\mathcal{D}_p]=m_if_t[\mathcal{D}_p,d_i]$, so the derivative of an output with respect to a local decay rate is one more appended node (verified numerically in Appendix~\ref{app:verification}). And memoization interpolates between walks and state vectors: storing the forward response at the states of a small ball around the actuator and enumerating walks from the sensor only until they reach that ball would trade memory for time continuously, at the price of a convolution on the time grid that the present exact scheme avoids.
